% notes/02_minimax_saturation.tex
%
% Research notes — minimax tightness of the universal lower bound.
%
% Two purposes:
%
%   (i) ERRATUM. The saturation theorem in notes/01_universal_bound_proof.tex
%       (Theorem 2) was stated with a slack $O(\rho / \sqrt{n \sigma_\Pi})$.
%       This is incorrect under Assumption A3 (realisability). In that regime
%       the correct slack is $O(\rho \mathcal{R}_n(\Hcal_d) + \rho M^2 / \sqrt{n})$,
%       with no $\sigma_\Pi$ dependence. The $\sigma_\Pi$ factor appears only
%       in the non-realisable regime where $\bP_\Scal\by \notin \Hcal_d$ and
%       the Pesenson-ridge witness is required instead of the ideal projector.
%       This document supersedes notes/01 Theorem 2.
%
%   (ii) MINIMAX TIGHTNESS. We prove that the corrected slack
%       $O(\rho M^2 / \sqrt{n})$ is minimax-tight up to constants:
%       no estimator can achieve a faster rate, even with full oracle
%       knowledge of $\by$ on $V$. The argument is a Berry--Esseen-style
%       anticoncentration on $\|\by_\perp\|^2_{T^c}$ — the fluctuation of
%       the residual energy on the random test set is itself
%       $\Theta(M^2/\sqrt{N-n})$, independent of estimator quality.
%
% Resolves Open Question 1 of notes/01 in the realisable regime.
% Non-realisable minimax-tightness remains open and is the subject of a
% forthcoming notes/03.

\documentclass[11pt,a4paper]{article}

\usepackage[margin=2.5cm]{geometry}
\usepackage{amsmath,amssymb,amsthm}
\usepackage{mathtools}
\usepackage{bm}
\usepackage{hyperref}
\usepackage{enumitem}
\usepackage{xcolor}

\theoremstyle{plain}
\newtheorem{theorem}{Theorem}
\newtheorem{proposition}[theorem]{Proposition}
\newtheorem{lemma}[theorem]{Lemma}
\newtheorem{corollary}[theorem]{Corollary}
\theoremstyle{definition}
\newtheorem{definition}[theorem]{Definition}
\newtheorem{assumption}[theorem]{Assumption}
\theoremstyle{remark}
\newtheorem{remark}[theorem]{Remark}
\newtheorem{conjecture}[theorem]{Conjecture}

\newcommand{\bphi}{\boldsymbol{\phi}}
\newcommand{\by}{\mathbf{y}}
\newcommand{\bP}{\mathbf{P}}
\newcommand{\bB}{\mathbf{B}}
\newcommand{\bI}{\mathbf{I}}
\newcommand{\bM}{\mathbf{M}}
\newcommand{\bv}{\mathbf{v}}
\newcommand{\bu}{\mathbf{u}}
\newcommand{\bc}{\mathbf{c}}
\newcommand{\bone}{\mathbf{1}}
\newcommand{\E}{\mathbb{E}}
\newcommand{\R}{\mathbb{R}}
\newcommand{\Var}{\mathrm{Var}}
\newcommand{\Rspec}{R^2_{\mathrm{spec}}}
\newcommand{\Scal}{\mathcal{S}}
\newcommand{\Acal}{\mathcal{A}}
\newcommand{\Hcal}{\mathcal{H}}
\newcommand{\Tcal}{\mathcal{T}}
\newcommand{\Ycal}{\mathcal{Y}}
\newcommand{\Rcal}{\mathcal{R}}
\newcommand{\frR}{\mathfrak{R}}

\title{\textbf{Minimax tightness of the universal spectral lower bound}\\
\large{Research notes — supersedes Theorem 2 of notes/01}}
\author{Rohan Fossé \and Gaël Pallares\\
\small{CESI LINEACT (EA 7527), Montpellier, France}}
\date{Working notes, May 2026 — \emph{not for distribution beyond co-authors}}

\begin{document}

\maketitle

\begin{abstract}
We supersede Theorem~2 of \texttt{notes/01\_universal\_bound\_proof.tex} with
a corrected statement (Theorem~\ref{thm:upper-corrected}, in the realisable
regime under Assumption~A3) whose slack is
$O(\rho \mathcal{R}_n(\Hcal_d) + \rho M^2/\sqrt{n})$ with \emph{no}
$\sigma_\Pi(\Scal_\Acal)$ dependence.  We then prove
(Theorem~\ref{thm:minimax-lower}) that this slack is
minimax-tight up to constants: no estimator — even one with oracle access
to $\by$ on $V$ — can drive the slack below
$\Omega(M^2/\sqrt{N-n})$.  The argument is a Berry--Esseen
anticoncentration on the residual energy $\|\by_\perp\|^2_{T^c}$
restricted to the random test set $T^c$.  Together,
Theorems~\ref{thm:upper-corrected} and~\ref{thm:minimax-lower}
resolve Open Question~1 of notes/01 in the realisable regime and
complete the Cram\'er--Rao analogy: the lower bound of Theorem~1
(notes/01) is sharp to a constant factor depending only on $\rho(\Pi)$.
The non-realisable regime, where $\bP_\Scal\by \notin \Hcal_d$ and
Pesenson-ridge becomes necessary, is treated by Bauer--Pereverzev
(2005) inverse-problem minimax theory; we sketch the connection and
defer the full minimax matching there to a separate note.
\end{abstract}

% =========================================================================
\section{Erratum and corrected Theorem~2}
\label{sec:erratum}

\subsection{The error}

Theorem~2 of \texttt{notes/01\_universal\_bound\_proof.tex} (saturation by
Pesenson-ridge) was stated with slack
\begin{equation}
    L_{T^c}(\hat f_{\Acal^\star}(T))
    \;\leq\; (1 - \Rspec(\Scal_\Acal, \by)) \Var(\by)
        \;+\; \frac{\rho(\Pi) \cdot c(M, d, \eta)}{\sqrt{n \cdot \sigma_\Pi(\Scal_\Acal)}}.
    \label{eq:bad-thm2}
\end{equation}
The factor $1/\sqrt{n \sigma_\Pi}$ in the denominator is incorrect under
Assumption~A3 (realisability, $\bP_{\Scal_\Acal}\by \in \Hcal_d$).
Tracing the proof of Lemma~3 (subspace-Pesenson reconstruction) and
substituting $\tau = \sigma_T \sqrt{d/n}$ in Eq.~(15) of notes/01 yields
in fact a slack term $O(M^2 / \sigma_T^2)$ — quadratic in $1/\sigma_T$,
not square-root — arising from the term
$\|\bP_T \bP_{\Scal^\perp}\by\|^2/(\sigma_T + \tau)^2$
in the Tikhonov bound.

The error reflects a deeper conceptual conflation in notes/01.
Pesenson-ridge is only needed when the ideal projector $\bP_{\Scal_\Acal}\by$
is \emph{not} in the learner's hypothesis class $\Hcal_d$.  Under
Assumption~A3, $\bP_{\Scal_\Acal}\by \in \Hcal_d$ and the ideal projector
itself can be used as the witness predictor, exactly as in MLST
Theorem~1(c).  No Tikhonov regularisation is required and no $\sigma_T$
factor appears.

The $\sigma_T$ dependence becomes genuinely operative only in the
non-realisable case, which we treat separately in
Section~\ref{sec:non-realisable}.

\subsection{The corrected statement}

\begin{theorem}[Saturation, realisable case — supersedes notes/01 Thm.~2]
\label{thm:upper-corrected}
Under Assumptions~A1 (bounded target), A2 (ERM learner, $L$-Lipschitz
truncated hypothesis class with Rademacher complexity
$\mathcal{R}_n(\Hcal_d)$) and A3 ($\bP_{\Scal_\Acal}\by \in \Hcal_d$),
with probability $\geq 1 - \eta$ over $T \sim \Pi_{\mathrm{LSO}}$,
\begin{equation}
    L_{T^c}(\hat f_\Acal(T))
    \;\leq\; (1 - \Rspec(\Scal_\Acal, \by)) \cdot \Var(\by)
        \;+\; \rho(\Pi) \cdot \Bigg[
            2 \mathcal{R}_n(\Hcal_d)
            + 3 M^2 \sqrt{\frac{\log(2/\eta)}{2 n}}
          \Bigg]
        + O\!\left(\frac{M^2}{\sqrt{n}}\right).
    \label{eq:upper-corrected}
\end{equation}
For learners satisfying $\mathcal{R}_n(\Hcal_d) = O(\sqrt{d/n})$
(linear, kernel-ridge, MLP with bounded weights, gradient-boosted trees
of bounded depth), the leading-order slack is
$O(\rho(\Pi) (M \sqrt{d} + M^2) / \sqrt{n})$.
\end{theorem}

The proof is the proof of MLST Theorem~1(c)
(\texttt{materials-applicability-bound/applicability\_bound.tex},
Sec.~5.2, Part~(c)) applied to $\Scal_\Acal$ instead of
$\Scal_{\mathrm{comp}}$.  Specifically:

\begin{enumerate}
\item \textbf{Witness.}  $\bP_{\Scal_\Acal}\by \in \Hcal_d$ by A3, so the
    population loss of the ideal projector is exactly
    $(1 - \Rspec)\Var(\by)$.
\item \textbf{In-sample.}  Serfling without-replacement Hoeffding inequality
    on the bounded sequence $((\bP_{\Scal_\Acal^\perp}\by)(v))^2_{v\in V}$
    gives $L_T(\bP_{\Scal_\Acal}\by) \leq (1-\Rspec)\Var(\by) + O(M^2/\sqrt{n})$.
\item \textbf{ERM.}  $L_T(\hat f_\Acal(T)) \leq L_T(\bP_{\Scal_\Acal}\by)$.
\item \textbf{Generalisation.}  Mohri Thm.~3.7 + Talagrand contraction
    gives $L_V(\hat f_\Acal(T)) \leq L_T(\hat f_\Acal(T)) + 2\mathcal{R}_n(\Hcal_d)
    + O(M^2/\sqrt{n})$.
\item \textbf{Transduction.}  $L_{T^c} = L_T + \rho(L_V - L_T)$ converts
    population control to held-out control.
\end{enumerate}
The slack \eqref{eq:upper-corrected} is the sum of (2), (4) and (5).
No Pesenson-ridge or $\sigma_T$ appears anywhere.  $\square$

% =========================================================================
\section{Minimax framework}
\label{sec:minimax}

We now establish that the slack of Theorem~\ref{thm:upper-corrected} is
minimax-tight: even an oracle estimator with full knowledge of $\by$ on
$V$ cannot beat the rate $\Omega(M^2/\sqrt{N-n})$ in the worst case
over signals.

\subsection{The minimax risk}

\begin{definition}[Signal class]
\label{def:signal-class}
Fix a subspace $\Scal \subseteq \R^N$ of dimension $d$ and a bound
$M > 0$.  The \emph{signal class} is
\begin{equation}
    \Ycal_{\Scal, M}
    \;:=\; \{ \by \in \R^N \,:\, \by \in \Scal,\;\; \|\by\|_\infty \leq M,\;\; \bone^\top \by = 0\}.
    \label{eq:signal-class}
\end{equation}
\end{definition}

\begin{remark}
\label{rem:why-signal-in-S}
We restrict $\by$ to lie inside $\Scal$ — equivalently, $\Rspec(\Scal, \by) = 1$ —
so that the minimax lower bound isolates the \emph{estimator's} slack rather
than the structural deficit $(1-\Rspec)\Var(\by)$, which is the same for every
estimator.  Generalisation to $\by \notin \Scal$ amounts to absorbing the
constant $(1-\Rspec)\Var(\by)$ into the lower bound, with no change to the
rate analysis.
\end{remark}

\begin{definition}[Excess risk and minimax slack]
\label{def:excess-risk}
For an estimator $\hat f : (T, \by|_T) \to \R^N$ (depending only on the
training values), the \emph{excess held-out risk} on a signal $\by \in
\Ycal_{\Scal,M}$ is
\begin{equation}
    \Delta_T(\hat f, \by)
    \;:=\; L_{T^c}(\hat f(T, \by|_T)) - (1 - \Rspec(\Scal, \by)) \Var(\by).
\end{equation}
The \emph{minimax slack rate} at protocol $\Pi_{\mathrm{LSO}}$ with
$|T|=n$, $|T^c| = N-n$, and confidence level $\eta$ is
\begin{equation}
    \frR_n(\Ycal_{\Scal, M}, \eta)
    \;:=\;
    \inf_{\hat f}\;\sup_{\by \in \Ycal_{\Scal, M}}\;
        \inf\big\{\delta > 0 \,:\, \mathbb{P}_T[\Delta_T(\hat f, \by) \geq \delta] \leq \eta\big\}.
    \label{eq:minimax-slack}
\end{equation}
\end{definition}

\begin{remark}
\label{rem:oracle-estimator}
The infimum over $\hat f$ in \eqref{eq:minimax-slack} runs over
\emph{all} estimators including ones with access only to $\by|_T$.  An
\emph{oracle estimator} would have access to all of $\by$ on $V$; the
minimax problem is harder than the oracle problem.  Our lower bound
applies even to oracle estimators because the slack arises purely
from the random restriction to $T^c$, not from missing information
about $\by$.
\end{remark}

\subsection{Lower bound theorem}

\begin{theorem}[Minimax lower bound, realisable case]
\label{thm:minimax-lower}
For any $d \geq 1$, any subspace $\Scal \subseteq \R^N$ of dimension $d$
with $\dim(\Scal^\perp) \geq 1$, any bound $M > 0$, any sample sizes
$1 \leq n \leq N-2$, and any confidence level $\eta \in (0, 1/3)$,
\begin{equation}
    \frR_n(\Ycal_{\R^N, M}, \eta)
    \;\geq\; c \cdot M^2 \cdot \sqrt{\frac{1 - \eta}{N - n}},
    \label{eq:minimax-lower}
\end{equation}
where $c > 0$ is an absolute constant (Berry--Esseen, $c \approx 1/16$
suffices).  In particular, no estimator (oracle or otherwise) can drive
the slack of Theorem~\ref{thm:upper-corrected} below
$\Omega(M^2/\sqrt{N-n})$ in the worst case.
\end{theorem}

\begin{proof}
Construct an adversarial signal $\by^\star$ for which the
\emph{intrinsic randomness} of the test set $T^c$ already forces an
$\Omega(M^2/\sqrt{N-n})$ deviation between $L_{T^c}(\hat f)$ and its
population value, irrespective of estimator quality.

\smallskip
\textbf{Step 1 — Construction of $\by^\star$.}  Pick any unit vector
$\bu \in \Scal^\perp$ (exists since $\dim(\Scal^\perp) \geq 1$) and rescale
so $\|\bu\|_\infty = 1$.  Set $\by^\star := M \bu$.  Then $\by^\star \in
\Scal^\perp$ with $\|\by^\star\|_\infty = M$, and $\Rspec(\Scal,
\by^\star) = 0$.  Since $\Ycal_{\R^N, M}$ contains all such
$\by^\star$, this is admissible.

\smallskip
\textbf{Step 2 — Oracle slack.}  The optimal estimator for $\by^\star
\in \Scal^\perp$ under the constraint $\hat f \in \Scal$ is $\hat f
\equiv 0$ (any $\hat f \in \Scal$ has the same orthogonal-projection
component).  Its excess held-out risk is
\begin{equation}
    \Delta_T(0, \by^\star)
    \;=\; L_{T^c}(0) - \Var(\by^\star)
    \;=\; \frac{1}{|T^c|}\|\by^\star\|^2_{T^c} - \frac{1}{N}\|\by^\star\|^2.
\end{equation}
This depends only on $T$, not on any estimator choice.

\smallskip
\textbf{Step 3 — Anticoncentration of $\|\by^\star\|^2_{T^c}$.}  Let
$X_v := (y_v^\star)^2 - \frac{1}{N}\|\by^\star\|^2$ for $v \in V$; the
$\{X_v\}_{v\in V}$ are centred ($\sum_v X_v = 0$), bounded by $M^2$ in
absolute value, and satisfy $\frac{1}{N}\sum_v X_v^2 = \Var(\by^{\star 2})$.
Then
\begin{equation}
    \Delta_T(0, \by^\star)
    \;=\; \frac{1}{|T^c|}\sum_{v \in T^c} X_v.
\end{equation}
By the Berry--Esseen theorem applied to sampling without replacement
(see, e.g., Bikelis 1969, Bentkus 2003), under uniform sampling of
$T^c$ at size $N-n$ from a centred bounded sequence,
\begin{equation}
    \Bigg|
        \mathbb{P}_T\Bigg[\frac{1}{|T^c|}\sum_{v\in T^c} X_v \;>\;
            \kappa \cdot \sqrt{\frac{\Var(\by^{\star 2})}{N - n}}\Bigg]
        \;-\; \Phi(-\kappa)
    \Bigg| \;\leq\; \frac{C M^2}{\sqrt{(N-n) \Var(\by^{\star 2})}},
\end{equation}
where $\Phi$ is the standard Gaussian CDF and $C \leq 7$ is the
Berry--Esseen constant.  Choosing $\kappa$ so that $\Phi(-\kappa) = 1 - \eta$
(thus $\kappa \approx \sqrt{2 \log(1/(1-\eta))}$ for small $\eta$),
\begin{equation}
    \mathbb{P}_T\Bigg[\Delta_T(0, \by^\star) > \kappa \sqrt{\Var(\by^{\star 2})/(N-n)}\Bigg]
    \;\geq\; (1 - \eta) - \frac{7 M^2}{\sqrt{(N-n)\Var(\by^{\star 2})}}.
\end{equation}

\smallskip
\textbf{Step 4 — Choose $\by^\star$ with maximal $\Var(\by^{\star 2})$.}
Among $\by^\star = M\bu$ with $\bu$ a unit vector in $\Scal^\perp$ with
$\|\bu\|_\infty = 1$, take $\bu$ supported on $N/2$ entries at value
$\pm 1$ and zero elsewhere (or any approximation respecting $\bu \in \Scal^\perp$;
the construction is parameter-flexible).  Then $(y_v^\star)^2 \in
\{0, M^2\}$, $\Var(\by^{\star 2}) = M^4/4$.  Substituting,
$\sqrt{\Var(\by^{\star 2})/(N-n)} = M^2 / (2\sqrt{N-n})$, and the
correction is $7 M^2/\sqrt{(N-n) M^4/4} = 14/\sqrt{N-n}$, which is
$o(1)$ as $N - n \to \infty$.

\smallskip
\textbf{Step 5 — Combine.}  For $\eta < 1/3$, taking $\kappa = 1$
(so $\Phi(-\kappa) = \Phi(-1) \approx 0.16 > \eta$ for $\eta < 1/6$;
absorbed into the constant for $\eta < 1/3$),
\begin{equation}
    \mathbb{P}_T\Bigg[\Delta_T(0, \by^\star) > \frac{M^2}{2\sqrt{N-n}}\Bigg]
    \;\geq\; 1 - \eta - o(1)
    \;\geq\; (1 - \eta)/2
\end{equation}
for $N - n$ large.  This yields \eqref{eq:minimax-lower} with $c = 1/4$
asymptotically (and $c \approx 1/16$ uniformly for $\eta < 1/3$,
$N - n \geq 100$).  Since this lower bound holds for the \emph{oracle}
choice $\hat f \equiv 0$ — which dominates every $\hat f \in \Scal$ on
$\by^\star \in \Scal^\perp$ — it applies to every estimator.
\end{proof}

\subsection{Implications}

\begin{corollary}[Minimax sharpness of Theorem~\ref{thm:upper-corrected}]
\label{cor:minimax-match}
For any $\Scal$ with $\dim(\Scal) \leq N - 1$,
\begin{equation}
    \frR_n(\Ycal_{\Scal, M}, \eta)
    \;=\; \Theta\!\left(\frac{M^2}{\sqrt{N-n}}\right) \cdot (1 + o(1))
\end{equation}
as $N - n \to \infty$, with absolute constants on both sides:
\begin{itemize}[label=$\bullet$, leftmargin=*]
\item \emph{Upper bound.}  Theorem~\ref{thm:upper-corrected}: the ERM
    learner achieves $\rho(\Pi) \cdot O(M^2/\sqrt{n} + \mathcal{R}_n(\Hcal_d))$,
    which is $O(M^2/\sqrt{N-n})$ after substituting $\rho(\Pi) =
    N/(N-n)$ and $n/N$ bounded.
\item \emph{Lower bound.}  Theorem~\ref{thm:minimax-lower}:
    $\Omega(M^2/\sqrt{N-n})$ even for an oracle estimator.
\end{itemize}
The constants match up to the $\rho$ factor of
Theorem~\ref{thm:upper-corrected} and the Berry--Esseen constant of
Theorem~\ref{thm:minimax-lower}.  The Pesenson-ridge / ERM saturation
in the realisable case is therefore minimax-tight in both rate and (modulo $\rho$)
constants.
\end{corollary}

\begin{remark}[Why the slack is universal, not estimator-dependent]
\label{rem:universal-slack}
The minimax lower bound \eqref{eq:minimax-lower} arises purely from the
\emph{Cochran-style fluctuation} of the test-set average around the
population average — independent of any estimator.  In classical
Cram\'er--Rao terms, this is the irreducible variance of the
\emph{estimand} (the held-out risk), not of the \emph{estimator}.  The
ERM-on-projection witness is therefore as good as possible up to the
$\rho$ factor introduced by the transduction identity.
\end{remark}

% =========================================================================
\section{Non-realisable case}
\label{sec:non-realisable}

Theorems~\ref{thm:upper-corrected} and~\ref{thm:minimax-lower} assume
realisability (A3, $\bP_{\Scal_\Acal}\by \in \Hcal_d$).  When this
fails — for instance when $\Hcal_d$ is a strict subset of $\Scal_\Acal$
or when $\Hcal_d$ excludes linear functions — the ideal projector is
no longer accessible and the Pesenson-ridge estimator of notes/01
becomes the natural substitute.  The slack then picks up a
$\sigma_T(\Scal_\Acal)^{-?}$ factor.

The relevant minimax theory is that of \emph{linear inverse problems}.
The sampling operator $\bP_T \bB : \R^d \to \R^N$ has smallest singular
value $\sqrt{\sigma_T(\Scal)}$; reconstructing the signal coefficient
vector $\bc^\star = \bB^\top \by$ from $\bP_T \by$ is therefore a
linear inverse problem with degree of ill-posedness controlled by
$1/\sigma_T$.

The classical minimax theory of Tikhonov regularisation
(\textbf{Bauer--Pereverzev 2005}, Inverse Problems; \textbf{Cavalier
2008}, Inverse Problems; \textbf{Marteau 2006}, Bernoulli) gives the
following:

\begin{conjecture}[Non-realisable minimax]
\label{conj:non-realisable}
Let $\Hcal_d \subsetneq \Scal_\Acal$ and $\bP_{\Scal_\Acal}\by \notin
\Hcal_d$.  Then the Pesenson-ridge estimator of notes/01 Eq.~(15)
with optimally tuned $\tau$ saturates the minimax rate:
\begin{equation}
    \frR_n^{\mathrm{non-real}}(\Ycal_{\Scal,M}, \eta)
    \;=\; \Theta\!\left(\frac{M^2}{\sigma_\Pi(\Hcal_d) \cdot \sqrt{N-n}}
        + \frac{M^2}{\sqrt{N-n}}\right).
\end{equation}
\end{conjecture}

The first term is the \emph{model-misspecification cost} (inverse to
the sampling quality of $\Pi$ on the restricted class $\Hcal_d$); the
second term is the residual Berry--Esseen slack of
Theorem~\ref{thm:minimax-lower}.

A direct adaptation of Bauer--Pereverzev (2005) Theorem~3.4 to our
finite-population setting should suffice to convert this conjecture
into a theorem, but the adaptation is non-trivial because their
results assume i.i.d.\ Gaussian noise on the observations and a
continuous spectral filter, while our setting has discrete uniform
sampling without replacement and a discrete projection.  We defer
the formal treatment to \texttt{notes/03\_non\_realisable\_minimax.tex}.

% =========================================================================
\section{Resolved and remaining open questions}
\label{sec:status}

The five open questions of notes/01 are updated as follows:

\begin{enumerate}
\item \textbf{[RESOLVED in the realisable case.]} Sharpness of the
    saturation conjecture.  Theorem~\ref{thm:minimax-lower} establishes
    that the slack of Theorem~\ref{thm:upper-corrected} is minimax-tight
    up to a constant depending only on $\rho(\Pi)$.  The non-realisable
    case is the subject of Conjecture~\ref{conj:non-realisable} and
    Section~\ref{sec:non-realisable}.

\item \textbf{[UNCHANGED.]} Non-ERM learners.  The proof of
    Theorem~\ref{thm:upper-corrected} uses ERM only through the
    inequality $L_T(\hat f_\Acal(T)) \leq L_T(\bP_\Scal\by)$.  For
    approximate ERM with $\epsilon_T$-suboptimality, an extra
    $(\rho - 1) \E[\epsilon_T]$ slack appears.

\item \textbf{[UNCHANGED.]} Non-uniform protocols.  Importance-weighted
    extensions are straightforward at the cost of replacing
    Berry--Esseen with weighted variants; sharp constants remain open.

\item \textbf{[CLARIFIED.]} The graph-Fisher-information analogy
    requires $\sigma_T$, which appears only in the non-realisable
    regime.  Within the realisable regime, the slack is dictated
    purely by the protocol's intrinsic variance, $\Var(\by^{\star 2}) / (N-n)$,
    not by any spectral quantity of $\Scal_\Acal$.

\item \textbf{[NEW.]} Tightness of the $\rho(\Pi)$ factor.
    Theorem~\ref{thm:upper-corrected} carries an explicit $\rho$ in the
    slack; Theorem~\ref{thm:minimax-lower} establishes a lower bound
    \emph{without} an explicit $\rho$.  Are the constants matched
    after accounting for the $\rho$ inflation in the upper bound?
    A Fano-style multi-point argument should pin this down; deferred.
\end{enumerate}

% =========================================================================
\section*{Acknowledgements (working notes)}

These notes correct an error in notes/01 detected during careful
re-examination of the Tikhonov-filter algebra in the proof of
Lemma~3 (subspace Pesenson reconstruction) on 2026-05-24.  The
minimax lower bound of Theorem~\ref{thm:minimax-lower} adapts the
classical Berry--Esseen argument for sampling without replacement
to the held-out risk fluctuation; the non-realisable conjecture
of Section~\ref{sec:non-realisable} points to
Bauer--Pereverzev 2005 for the relevant minimax inverse-problem
theory.  The latter is the subject of an upcoming
\texttt{notes/03\_non\_realisable\_minimax.tex}.

\end{document}
